\documentclass[conference]{IEEEtran}
	\usepackage[utf8]{inputenc}
	\usepackage[T1]{fontenc}
	\usepackage{graphicx}
	\usepackage{cite}
	\usepackage{amsmath, amssymb}
	\usepackage{booktabs}
	\usepackage{tikz}
	\usepackage{pgfplots}
	\usepackage{subcaption}
	\usepackage{xcolor}
	\usepackage{enumitem}
	\usepackage{pgfplotstable}
	\pgfplotsset{compat=1.17}
	\usetikzlibrary{positioning, arrows.meta, shapes.geometric, shapes.multipart, calc}
	\usepackage{float}
	\usepackage{algorithm}
	\usepackage{algorithmic}
	\renewcommand{\algorithmicrequire}{\textbf{Input:}}
	\renewcommand{\algorithmicensure}{\textbf{Output:}}
	\usepackage{pifont}
	% --- Hyperref (load last in preamble) ---
	\usepackage[pdfencoding=auto,unicode,psdextra,colorlinks=true, linkcolor=blue, citecolor=blue, urlcolor=blue]{hyperref}
	
	% --- TITLE ---
	\title{FedSemGNN: A Scalable, Semantic-Aware Federated RL Framework for Efficient 6G Edge Orchestration}
	
	
	% --- AUTHORS ---
	\author{
		Qaiser Muhammad Abdur Rehman\textsuperscript{a}, 
		Mingchu Li\textsuperscript{a,b*}, 
		Sanam Shahla Rizvi\textsuperscript{c}, 
		and Se Jin Kwon\textsuperscript{d} \\[4pt]
		\textsuperscript{a}School of Software, Dalian University of Technology, Dalian 116024, China\\
		\textsuperscript{b}School of Computer and Information Engineering, Jiangxi Normal University, Nanchang, Jiangxi, China\\
		\textsuperscript{c}Raptor Interactive (Pty) Ltd., Centurion 0157, South Africa\\
		\textsuperscript{d}Department of Computer Science \& Artificial Intelligence, Dongguk University, Seoul 04620, South Korea\\[6pt]
		\textit{Emails:} qaisermuhammad@hotmail.com, mingchu@jxnu.edu.cn, 
		sanam\_shahla@hotmail.com, skwon1109@gmail.com\\[2pt]
		*Corresponding author: mingchu@jxnu.edu.cn
	}
	
	
	
	\begin{document}
		
		\maketitle
		
		\begin{abstract}
			6G edge orchestration must place heterogeneous services under strict latency, energy, and privacy constraints at massive device density. We present FedSemGNN, a semantic-aware hierarchical federated reinforcement learning framework that integrates continual semantic task embeddings, graph convolutional topology encoding, and two-level federated PPO synchronization to reduce communication overhead while preserving convergence stability.
			
			FedSemGNN exposes tunable semantic matching ($\tau$) and continual-learning regularization ($\lambda_{\mathrm{EWC}}$) and supports priority-aware scheduling via latency weighting and priority-adaptive semantic thresholds for URLLC-style services. Implemented in EdgeSimPy, FedSemGNN is evaluated over 1,000 orchestration timesteps and compared with FlatFedPPO, HierFedPPO, HSQF, RandomPlacement, and a centralized PPO baseline. Results show improved reward and semantic fidelity with sub-millisecond decision latency, reduced migrations, and low communication footprint while scaling to 10,000 nodes with near-constant control-plane latency.
		\end{abstract}
		
		
		\begin{IEEEkeywords}
			Federated Reinforcement Learning, Semantic Task Orchestration, Graph Neural Networks, Edge Intelligence, Proximal Policy Optimization, 6G Networks
		\end{IEEEkeywords}
		
		
		\section{Introduction}
		
		The imminent arrival of 6G networks promises to enable next-generation edge intelligence with ultra-low latency, high reliability, and context-aware services including extended reality, autonomous systems, and tactile internet applications~\cite{cui2025}. These use cases demand fine-grained orchestration of computing and data services across heterogeneous edge nodes~\cite{shi2023}. Traditional cloud-based approaches cannot meet such requirements due to latency bottlenecks, reinforcing the importance of edge computing to bring processing closer to end users~\cite{vahabi2025}.
		
		The dynamic and resource-constrained nature of edge environments combined with energy constraints, privacy requirements, and workload variability introduces substantial orchestration challenges~\cite{yang2023rlcs}. Recent studies have investigated reinforcement learning for adaptive orchestration, allowing real-time learning of placement and migration strategies~\cite{shi2023}. However, centralized RL methods suffer from scalability limitations and cannot preserve user privacy. Furthermore, most RL-based orchestration approaches disregard semantic content in service tasks, leading to suboptimal placements in latency-sensitive or context-dependent scenarios such as extended reality or real-time analytics~\cite{zhang2025semantic}.
		
		Federated reinforcement learning provides a promising alternative by distributing policy learning across edge nodes without sharing raw data~\cite{wu2024topology}. While FRL frameworks have shown promise in handling privacy and scalability, they frequently overlook two critical dimensions: semantic task awareness and network topology encoding~\cite{ogenyi2025comprehensive, zhang2024optimizing}. Modern edge workloads often include semantic information such as task intent or content category, which if properly modeled can significantly enhance placement fidelity and decision stability~\cite{zhang2025semantic}. At the same time, orchestration effectiveness depends on network structure, including link quality, proximity, and cluster connectivity. Yet most FRL methods treat edge nodes as unstructured agents without considering topological relationships~\cite{wu2024topology}.
		
		Graph neural networks offer an effective means to encode network topology, but their integration into federated, real-time orchestration remains limited~\cite{makris2024gnnrl, zhang2024optimizing}. Moreover, frequent global synchronization in federated learning systems often results in high communication overhead and latency. Hierarchical FL frameworks reduce this burden by performing local aggregation within clusters before periodic global updates~\cite{sad2024fearless, he2022acefl}. However, few approaches combine hierarchical control with semantic embeddings and graph-based topology encoding in a unified orchestration framework~\cite{wu2024topology, vahabi2025}.
		
		Existing studies also suffer from limited experimental scope. Many evaluations use synthetic environments with minimal orchestration realism or baseline diversity. Large-scale scenarios are rarely validated beyond small-scale simulations, and systematic scalability analysis using mathematical extrapolation from empirically validated base experiments remains uncommon in the federated edge orchestration literature.~\cite{yang2023rlcs, shi2023}, 
		
		\noindent\textbf{Contributions.} We propose \textbf{FedSemGNN}, a federated reinforcement learning framework designed for semantic-aware, topology-informed, and hierarchical orchestration in 6G edge environments. Beyond the core architecture, we explicitly expose and analyze key hyperparameters (semantic matching threshold and continual-learning regularization) and incorporate task priority for 6G URLLC scenarios. The key contributions of this work include:
		
		\begin{itemize}
			\item \textbf{Semantic task embedding with continual learning:} Service requests are embedded into a compact semantic space (default dimension 16) using a continual learning encoder that extracts features from resource demands, service types, and temporal patterns. The encoder uses elastic weight consolidation (EWC) with regularization $\lambda_{\mathrm{EWC}}$ to suppress catastrophic forgetting while adapting to emerging workload characteristics; we report sensitivity to $\lambda_{\mathrm{EWC}}$.
			
			\item \textbf{Graph convolutional network topology encoding:} A two-layer GCN with symmetric normalization captures structural relationships including bandwidth, affinity, and proximity. The architecture processes node features through 128 to 64 to 64 dimensional transformations with ReLU activation, allowing placement decisions that consider network-wide constraints and multi-hop dependencies across mesh-connected edge infrastructure.
			
			\item \textbf{Hierarchical federated PPO with dual-level synchronization:} The framework employs two-level aggregation (default $K_1{=}10$, $K_2{=}50$) to reduce communication overhead while maintaining convergence stability. The design supports scalable federated learning without requiring fixed cluster assignments.
			
			\item \textbf{Priority-aware orchestration primitives:} We model service priority (emergency vs. ordinary) and incorporate priority-weighted latency in the reward. We additionally enable priority-adaptive semantic thresholds to better satisfy URLLC-like service constraints.
			
			\item \textbf{Comprehensive experimental validation:} We implement FedSemGNN in EdgeSimPy~\cite{souza2023edgesimpy}, a high-fidelity edge computing simulator, with detailed instrumentation of control-plane latency, reward trajectories, semantic fidelity, power consumption, communication footprint, and migration dynamics. The base evaluation uses a heterogeneous edge infrastructure with mobile users following realistic mobility patterns over 1,000 orchestration timesteps.
			
			\item \textbf{Exceptional near-constant scalability:} Scalability from 100 nodes to 10,000 nodes is validated through simulated experiments with dynamically generated topologies representing IoT and enterprise edge scenarios. FedSemGNN achieves near-constant orchestration latency (0.309--0.328 ms) with 0.967$\times$ scaling factor across 100$\times$ scale range, demonstrating sub-linear computational complexity. This shows that compact neural architectures and hierarchical aggregation enable massive-scale simulation performance. FedSemGNN is compared against four competitive baselines: FlatFedPPO, HierFedPPO, HSQF, and RandomPlacement.
			
			\item \textbf{Superior performance across all metrics:} FedSemGNN achieves normalized reward of 0.91 (8.4 percent higher than FlatFedPPO), orchestration latency of 0.36 ms per decision (315 times faster than FlatFedPPO), communication efficiency of 1.394 reward per MB (174 times higher than FlatFedPPO), 100 percent semantic fidelity with zero migrations (compared to 6,840 in FlatFedPPO), and energy consumption of 72.1 W (36.2 times lower than FlatFedPPO).
		\end{itemize}
		
		Unlike previous methods that treat semantics or topology heuristically, FedSemGNN unifies both within a scalable federated RL framework with mathematically grounded reward optimization. The reward function jointly minimizes power consumption and latency while encouraging semantic fidelity through state representation design. By embedding semantic structure into the learning process and encoding topological constraints through graph neural networks, FedSemGNN supports responsive orchestration with minimal latency and communication cost. The system demonstrates real-time operation under 0.5 ms per decision, offering a viable foundation for control-plane design in 6G edge computing.
		
		\noindent\textbf{Paper organization.} The remainder of this paper is structured as follows: Section~\ref{sec:related} reviews related work in federated reinforcement learning, semantic scheduling, and graph neural network based orchestration. Section~\ref{sec:system_architecture} presents the FedSemGNN system architecture including semantic embedding, GNN topology encoding, and hierarchical PPO coordination. Section~\ref{sec:methodology} details the proposed methodology with mathematical formulations for reward optimization, policy learning, and federated aggregation. Section~\ref{sec:evaluation} discusses empirical evaluations across baseline comparisons and scalability analysis. Section~\ref{sec:conclusion} concludes the paper, and Section~\ref{sec:limitations_future} outlines limitations and future research directions.
		
		
		\section{Related Work}
		\label{sec:related}
		
		\subsection{Federated Reinforcement Learning for Edge Orchestration}
		
		Federated reinforcement learning (FRL) has emerged as a robust approach for edge orchestration under privacy constraints. Unlike centralized RL, FRL enables distributed agents to train local policies without exposing raw data, which is critical for 6G-enabled services~\cite{kianpisheh2024,huang2022collaborative}. Recent studies explore FRL in hybrid cloud-edge environments~\cite{wang2021,consul2024hybrid} and under hierarchical model training setups~\cite{he2022acefl}, with applications including anomaly detection, task offloading, and digital twin optimization~\cite{lu2020low,prathiba2022}. Despite growing interest, most current methods still rely on abstract or flat feature vectors and struggle to encode spatial or semantic structure, limiting coordination under dynamic network topologies~\cite{yan2023v2v,sad2024fearless}.
		
		\subsection{Semantic-Aware Scheduling and Service Placement}
		
		Semantic embeddings offer robust generalization across orchestration tasks by encoding high-level intent and contextual similarity. Early works have applied semantic modeling in edge caching and service personalization~\cite{prathiba2022,tam2025graph}. Semantic compression has also been explored for communication-aware model training~\cite{wang2021}, but integration into reinforcement learning pipelines for distributed scheduling remains sparse. Recent advances suggest that semantic representations can reduce service migrations and enhance policy transfer across domains~\cite{zhang2025semantic,li2021network}. Our framework builds on this insight by coupling semantic representations with federated PPO for personalized orchestration policies.
		
		\subsection{GNN-Based Topology Encoding for Edge Intelligence}
		
		Graph neural networks (GNNs) are increasingly used to capture spatial relationships in wireless and edge environments. They support topology-aware orchestration by encoding node states, link quality, and global dependencies~\cite{zhang2024optimizing,jere2023distributed}. GNN-RL combinations have shown promise in multi-hop offloading, edge caching, and vehicular scheduling~\cite{makris2024gnnrl,hao2021urllc}, though few operate under federated constraints. Additionally, federated orchestration often lacks topology adaptation, limiting robustness under dynamic infrastructure changes. Our work integrates GNN-based topology modeling into a federated PPO loop, enabling communication-efficient and topology-consistent policy learning.
		
		
		\subsection{Hierarchical Federated Learning for Communication Efficiency}
		
		Hierarchical federated learning (HFL) reduces synchronization overhead by grouping clients into subregions that perform local model updates before periodic global aggregation. This model is especially suited for edge and IoT systems, where communication cost and energy efficiency are paramount and where clients exhibit heterogeneous energy, compute, and mobility characteristics~\cite{tan2023selfadaptive,huang2022workercentric}. While HFL has shown promise in supervised learning, its use in reinforcement learning for real-time orchestration remains limited~\cite{tam2025graph}. Existing methods often neglect dynamic service behaviors or lack mechanisms for topology-informed coordination. In our work, we extend HFL principles to federated PPO and apply it in a two-tier control structure with local synchronization every 10 epochs and global synchronization every 50 epochs. This results in significant communication savings, reducing total communication volume to 0.65 MB over 1,000 orchestration steps compared to 105.0 MB for flat federated PPO, representing a 161 times reduction in communication overhead.
		
		
		\subsection{Deep RL-Based Edge Offloading and Resource Allocation}
		
		Recent work on deep reinforcement learning for edge computing has demonstrated significant advances in task offloading and resource allocation. ECO-SDIoT~\cite{zhu2023eco} employs Double DQN with software-defined networking for global-view optimization, achieving 20--60 ms task offloading latency in IoT environments. While effective for batch-oriented offloading, ECO-SDIoT lacks semantic awareness and operates without federated privacy guarantees. GFL-LFF~\cite{regan2024gfl} combines GNN with federated learning for privacy-preserving data aggregation in Industrial IoT, reporting 1.6 s latency for secure data sharing. Although GFL-LFF demonstrates the benefits of GNN integration with federated approaches, it focuses on data aggregation rather than real-time orchestration. FRPVC~\cite{qian2025frpvc} applies federated deep reinforcement learning with DDQN for cost-efficient video caching in energy-constrained mobile edge networks, optimizing cache hit rates through denoised autoencoders and MDP-based resource allocation.
		
		While these methods demonstrate effective deep RL optimization for specific edge applications, they address narrow use cases (offloading, data aggregation, caching) without generalizing to semantic-aware task placement. Moreover, value-based methods (Double DQN, DDQN) can suffer from instability in high-dimensional action spaces compared to policy gradient approaches. FedSemGNN advances beyond these works by: (1) incorporating explicit semantic embeddings for application-aware orchestration, (2) employing hierarchical PPO for more stable policy learning than value-based methods, (3) using GNN topology encoding for efficient network representation, and (4) achieving sub-millisecond latency (0.36 ms) representing 55 to 4,444 times improvements over state-of-the-art methods while maintaining federated privacy guarantees.
		
		\noindent\textbf{Centralized RL and MARL context.} Reviewer feedback requested comparisons against mainstream centralized RL (e.g., SAC, TD3) and MARL scheduling. SAC~\cite{haarnoja2018sac} and TD3~\cite{fujimoto2018td3} are widely used for continuous-control problems, while MAPPO~\cite{yu2021mappo} is a strong multi-agent PPO variant. Our orchestration action is a large, combinatorial placement decision over heterogeneous edge servers, and our primary focus is privacy-preserving, topology-aware federation under strict control-plane latency. We therefore include a centralized PPO baseline as a directly comparable centralized policy-gradient control baseline (same observation construction and decision latency accounting) and discuss extensions to additional centralized RL and MARL baselines as future work.
		
		\noindent\textbf{Additional 6G-related works.} Beyond edge orchestration papers, 6G networking increasingly relies on intelligent optimization at the physical and network layers. For example, ANNS~\cite{chen2025anns} proposes an intelligent non-convex optimization scheme for IRS-aided next-generation mobile networks, highlighting the broader trend of learning-assisted optimization that motivates cross-layer intelligent orchestration.
		
		
		\subsection{Positioning of Our Work}
		
		To the best of our knowledge, no existing system jointly integrates semantic embeddings, GNN-based topology encoding, and hierarchical federated reinforcement learning within a unified orchestration framework. Our proposed \textbf{FedSemGNN} is the first to combine these dimensions in a federated setting, implemented on the EdgeSimPy simulator~\cite{souza2023edgesimpy}. The system supports detailed instrumentation across reward, latency, communication, power, and fidelity metrics. Evaluated over 1,000 orchestration steps on a heterogeneous edge infrastructure with scalability validated across 100 to 10,000 nodes (100$\times$ scale range) through simulated experiments, FedSemGNN demonstrates substantial improvements compared to both self-implemented baselines (FlatFedPPO, HierFedPPO, HSQF, RandomPlacement) and state-of-the-art published methods (ECO-SDIoT, GFL-LFF, FRPVC). Key simulation results include normalized reward of 0.91 (8.4 percent higher than FlatFedPPO), near-constant orchestration latency of 0.309--0.328 ms across all scales with 0.967$\times$ scaling factor (315 times faster than FlatFedPPO, 55 to 166 times faster than ECO-SDIoT, and 4,444 times faster than GFL-LFF), 100 percent semantic fidelity with zero migrations, communication efficiency of 1.394 reward per MB (76 times better than ECO-SDIoT), and energy consumption of 72.1 W (36.2 times lower than FlatFedPPO). These findings highlight the exceptional scalability and architectural potential of semantic and topological modeling in federated edge orchestration.
		
		
		\section{System Architecture}
		\label{sec:system_architecture}
		
		FedSemGNN enables scalable, low-latency, and privacy-preserving orchestration for 6G edge networks by integrating three core components:
		
		\begin{enumerate}
			\item \textbf{Semantic-Aware Task Representation}: Custom continual learning encoder extracts compact semantic embeddings from service characteristics to enable intent-aware orchestration~\cite{yan2023v2v, prathiba2022, li2021network},
			\item \textbf{GNN-Based Topology Encoding}: Graph convolutional network captures spatial dependencies and inter-node relationships across the edge infrastructure~\cite{makris2024gnnrl, tam2025graph}, and
			\item \textbf{Hierarchical Federated PPO}: Two-level synchronization coordinates distributed policy learning while preserving privacy and minimizing communication overhead~\cite{sad2024fearless, kianpisheh2024, yan2023v2v}.
		\end{enumerate}
		
		The architecture follows a hierarchical federated learning paradigm where multiple edge clusters operate semi-autonomously while maintaining global coordination through periodic aggregation. Each cluster contains resource monitors, edge nodes with local orchestrators, and graph-based topology encoders. User devices equipped with semantic encoders generate service requests that are processed through the orchestration pipeline.
		
		Local orchestrators make placement decisions using policies learned through proximal policy optimization, leveraging both semantic compatibility and topological context. The two-level synchronization structure—intra-cluster coordination and global aggregation—balances responsiveness with communication efficiency, enabling the system to scale to large simulated edge infrastructures while maintaining orchestration quality.
		
		\begin{figure*}[htbp]
			\centering
			\includegraphics[width=1.0\textwidth]{figures/architecture.png}
			\caption{Complete FedSemGNN system architecture showing hierarchical federated learning with semantic embeddings, GNN topology encoding, and two-level synchronization. User devices with semantic encoder feed multiple edge clusters (Cluster 1, Cluster 2), each containing resource monitors, edge nodes, GCN encoders, and local PPO agents. Intra-cluster sync occurs every $K_1=10$ epochs, while global DP aggregation and policy distribution occur every $K_2=50$ epochs.}
			\label{fig:system_framework}
		\end{figure*}
		
		\subsection{Semantic-Aware Task Representation}
		\label{subsec:semantic_representation}
		
		Service requests entering the FedSemGNN system are transformed into compact semantic embeddings that capture task intent and requirements. The semantic encoder processes heterogeneous service characteristics including resource demands (CPU, memory, network bandwidth), functional service types (web services, data analytics, real-time processing), temporal patterns (time-of-day preferences, periodic workloads), and quality-of-service requirements (latency sensitivity, throughput constraints).
		
		The encoder employs continual learning mechanisms to adapt to evolving service landscapes without forgetting previously learned patterns~\cite{yan2023v2v, prathiba2022}. When new service types emerge, the system preserves knowledge of existing services while incorporating novel patterns through regularization and experience replay. This adaptive capability is essential for dynamic edge environments where application requirements continuously evolve.
		
		Each edge node maintains a semantic capability profile that reflects its functional specialization and historical workload patterns. Semantic matching identifies compatible task-server pairs, enabling the system to maintain high placement fidelity by aligning service requirements with node capabilities~\cite{li2021network, jere2023distributed}. Tasks are considered semantically compatible with nodes when their embedding similarity exceeds a predefined threshold, ensuring quality-of-service consistency.
		
		Figure~\ref{fig:semantic_pipeline} illustrates the complete semantic processing pipeline from service arrival to placement decision. The pipeline begins with service/task input and graph/context features, which are processed by the semantic encoder to generate initial embeddings. These embeddings are then refined by the GNN encoder to incorporate topological context, followed by a projection head that produces the final 16-dimensional semantic vector $\mathbf{z}$. The semantic vector undergoes similarity search against server capability profiles to identify compatible placement candidates, culminating in the final placement decision that selects both the target cluster and specific edge node. A critical feedback loop (shown in purple dotted line) enables semantic adaptation: performance signals and policy outcomes are fed back to the semantic encoder, allowing the system to continuously refine its embedding space based on actual orchestration results. This adaptive mechanism is key to FedSemGNN's ability to maintain perfect semantic fidelity while adapting to evolving workload patterns.
		
		\begin{figure*}[htbp]
			\centering
			\includegraphics[width=0.95\textwidth]{figures/semantic_pipeline.png}
			\caption{Semantic embedding pipeline showing the flow from service requests through the custom continual learning encoder ($128 \rightarrow 64 \rightarrow 64 \rightarrow 16$), GNN encoding, projection head, semantic vector generation, similarity-based matching, and final placement decision. The feedback loop (purple dotted line) enables continuous semantic adaptation based on orchestration performance.}
			\label{fig:semantic_pipeline}
		\end{figure*}
		
		\subsection{Topology Encoding with Graph Neural Networks}
		\label{subsec:topology_encoding}
		
		The edge computing infrastructure is modeled as a graph where nodes represent compute servers and edges capture communication links~\cite{makris2024gnnrl, tam2025graph} with associated costs (latency, bandwidth, reliability). This graph representation enables the system to reason about spatial relationships, network proximity, and resource distribution across the infrastructure.
		
		Graph neural networks process this topological structure through iterative neighborhood aggregation, where each node accumulates information from its local neighborhood to build progressively richer representations~\cite{makris2024gnnrl, huang2022collaborative}. Initial node features include semantic affinity scores, current resource availability (CPU, memory, bandwidth), and historical usage patterns. Through multiple aggregation layers, nodes develop embeddings that encode both their local state and broader network context from multi-hop neighborhoods.
		
		The topology-aware embeddings enable spatially coherent orchestration decisions. By accounting for network proximity and communication costs, the system can place related services near each other and avoid bottleneck links. This spatial awareness is critical for reducing service migrations and maintaining stable placements even as workloads shift across the infrastructure.
		
		\subsection{Hierarchical Federated Proximal Policy Optimization}
		\label{subsec:hierarchical_ppo}
		
		The orchestration framework employs hierarchical federated learning to coordinate distributed policy learning across edge clusters. Each local orchestrator maintains its own policy for placement decisions, learning from local experiences while periodically synchronizing with neighboring orchestrators and the global coordinator.
		
		The hierarchical structure operates at two levels: intra-cluster coordination maintains coherence among orchestrators within the same edge cluster, while global aggregation ensures consistency across all clusters. This two-level design reduces communication overhead compared to flat federated architectures where all agents must synchronize with a central server for every update~\cite{sad2024fearless, yan2023v2v}. Local coordination enables rapid adaptation to cluster-specific workload patterns, while global aggregation prevents policy divergence and ensures system-wide optimization.
		
		Privacy preservation is achieved through secure aggregation protocols where only model parameters (not raw data or trajectories) are transmitted during synchronization. Differential privacy mechanisms add calibrated noise during aggregation to provide formal privacy guarantees. Gradient clipping prevents any single agent from dominating the aggregated policy, ensuring fairness and stability~\cite{mao2024optimizing, li2021network}.
		
		The policy learning algorithm employs proximal policy optimization, which constrains policy updates to prevent destructively large changes that could destabilize training~\cite{singh2024drl, kianpisheh2024}. Each orchestrator maintains an actor-critic architecture: the actor network selects placement actions based on current state observations, while the critic network estimates state values to guide policy improvement. The compact neural architecture enables low-latency decision-making suitable for real-time orchestration requirements.
		
		\subsection{Workflow Summary}
		\label{subsec:workflow}
		
		Figure~\ref{fig:workflow} presents the complete end-to-end orchestration workflow showing how system components interact during service placement. The workflow begins when a service request arrives at the system and proceeds through semantic encoding, topology-aware state fusion, graph neural network processing, and policy-based decision-making. After service placement, the system evaluates placement quality through reward computation, updates the local policy, and periodically synchronizes with other orchestrators.
		
		Two critical feedback mechanisms enable continuous system improvement. The global policy broadcast ensures all local orchestrators remain coordinated by distributing aggregated parameters from the central coordinator. The semantic adaptation feedback allows the embedding encoder to refine its representations based on actual orchestration outcomes, creating a closed-loop learning system that continuously improves semantic fidelity.
		
		Synchronization decisions occur at strategic intervals to balance responsiveness with communication efficiency. Intra-cluster synchronization maintains local coherence, while global aggregation ensures system-wide consistency without requiring constant communication. This adaptive synchronization strategy allows the simulated system to scale to large network sizes while maintaining orchestration quality~\cite{souza2023edgesimpy}.
		
		\begin{figure}[!t]
			\centering
			\includegraphics[width=0.50\textwidth]{figures/workflowdiagram.png}
			\caption{Complete FedSemGNN orchestration workflow showing the decision cycle: (0) Service Arrival $\rightarrow$ (1) Semantic Encoding $\rightarrow$ (2) GNN Encoding $\rightarrow$ (3) Cluster Selection $\rightarrow$ (4) Placement $\rightarrow$ (5) Reward Computation $\rightarrow$ (6) Local PPO Update $\rightarrow$ (7) Sync Decision $\rightarrow$ (7a-7b) Hierarchical Aggregation $\rightarrow$ (8) Policy Broadcast. Purple lines show critical feedback: solid for global policy broadcast to local agents, dotted for semantic adaptation feedback to encoder.}
			\label{fig:workflow}
		\end{figure}
		
		\subsection{Design Rationale and System Characteristics}
		\label{subsec:design_rationale}
		
		The FedSemGNN architecture is designed to address three fundamental challenges in 6G edge orchestration: semantic heterogeneity of service requests, dynamic network topology, and distributed coordination at scale.
		
		\textbf{Semantic Awareness:} Traditional resource-based orchestration treats all services homogeneously, considering only CPU, memory, and bandwidth requirements. FedSemGNN's semantic layer enables intent-aware placement by matching service characteristics with node capabilities, improving quality-of-service consistency and reducing placement violations~\cite{regan2024gfl, zhu2023eco}.
		
		\textbf{Topology Encoding:} Graph neural networks provide a principled approach to encode network structure, enabling the system to reason about spatial relationships and communication costs. This topological awareness reduces unnecessary service migrations by accounting for network proximity when making placement decisions.
		
		\textbf{Hierarchical Federation:} The two-level synchronization structure addresses the scalability-consistency tradeoff inherent in distributed learning. Frequent local coordination maintains cluster coherence, while periodic global aggregation ensures system-wide optimization without excessive communication overhead.
		
		\textbf{Compact Architecture:} The neural components employ small network architectures (16-64 hidden units) to enable sub-millisecond decision latency. This design choice prioritizes inference speed over model capacity, recognizing that orchestration decisions must be made in real-time as service requests arrive.
		
		\textbf{Privacy Preservation:} Secure aggregation and differential privacy mechanisms ensure that no agent can infer information about other agents' local workloads or placement decisions. This privacy guarantee would be essential for multi-tenant edge infrastructures where competing service providers share infrastructure.
		
		The combination of these design choices enables FedSemGNN to deliver scalable, low-latency orchestration while maintaining semantic fidelity, privacy, and energy efficiency. The architecture is modular, allowing individual components to be upgraded or replaced without redesigning the entire system.
		
		\subsection{Compatibility with Existing 6G Edge Platforms}
		FedSemGNN is designed as a control-plane decision module that consumes telemetry (resource utilization, link statistics, mobility context, and service descriptors) and outputs placement/migration actions. This interface aligns with common 6G edge stacks where monitoring and actuation are separated: telemetry can be collected from SDN controllers and edge orchestrators, while placements can be enacted by container orchestrators (e.g., Kubernetes) or MEC platforms. In deployment, the semantic encoder and policy inference can run in a lightweight controller pod; federated aggregation can be implemented as a periodic background job across sites (clusters) to respect privacy and bandwidth constraints.
		
		
		\section{Proposed Methodology}
		\label{sec:methodology}
		
		\subsection{System Model}
		We consider an edge computing network composed of $N$ heterogeneous edge servers interconnected by communication links with varying latency and bandwidth. Each server $v \in V$ is defined by a dynamic state vector $\mathbf{x}_v(t)$ that includes:
		\begin{itemize}
			\item CPU utilization and available processing capacity,
			\item Memory availability and active allocation,
			\item Uplink and downlink bandwidth,
			\item Historical workload statistics,
			\item Semantic service affinity scores.
		\end{itemize}
		
		The edge fabric is modeled as a weighted graph $G = (V, E)$ where edge weights encode communication costs (latency and bandwidth). To improve scalability and reduce synchronization overhead, the network employs hierarchical federated learning with two-level aggregation. Local orchestrators perform task placement before participating in periodic global policy aggregation, adapting dynamically to network topology and workload patterns~\cite{zhou2024qos, yurdem2024federated, zhou2024robustfl}.
		
		\subsection{Semantic Embedding}
		For each incoming task $i$, service characteristics are embedded into a compact semantic space with dimension $d_s$ (default $d_s=16$) using a continual learning encoder:
		\begin{equation}
			\mathbf{s}_i = f_{\text{embed}}(\text{task}_i) \in \mathbb{R}^{d_s},
		\end{equation}
		where $f_{\text{embed}}(\cdot)$ processes service features including resource demands, service type indicators, temporal features, and workload characteristics (including priority and latency sensitivity). The encoder employs elastic weight consolidation (EWC) with regularization parameter $\lambda_{\text{EWC}}$ (default $0.4$) to prevent catastrophic forgetting when adapting to new service types. These values are not fixed constants: we expose them as hyperparameters and report sensitivity analysis.
		
		Each edge node $v$ maintains a semantic profile vector $\mathbf{p}_v \in \mathbb{R}^{16}$ that reflects its service capabilities and historical workload characteristics. Semantic affinity between task $i$ and node $v$ is computed using cosine similarity:
		\begin{equation}
			\text{sim}(i,v) = \frac{\mathbf{s}_i \cdot \mathbf{p}_v}{\|\mathbf{s}_i\|_2 \|\mathbf{p}_v\|_2},
		\end{equation}
		where the dot product $\mathbf{s}_i \cdot \mathbf{p}_v$ measures alignment between task and node capabilities, normalized by their Euclidean norms to produce values in $[-1, 1]$. Values close to 1 indicate strong semantic match, while values near 0 suggest orthogonal service types. The semantic score is fused with numerical resource metrics to form a unified state representation supporting intent-aware orchestration~\cite{grasso2022hhome}.
		
		\subsubsection{Semantic similarity threshold and adaptation}
		To decide whether a placement satisfies semantic requirements, we use a similarity threshold $\tau$ such that a task--server pair is considered semantically compatible if $\mathrm{sim}(i,v)\ge\tau$. In this work, $\tau$ is treated as a tunable system knob (default $0.3$) rather than an arbitrary constant. We report sensitivity to $\tau$ and use it to trade off strict semantic matching (higher $\tau$) versus feasibility and flexibility (lower $\tau$).
		
		To reflect heterogeneous 6G services, we additionally support priority-adaptive thresholds. Each service is assigned a priority $p_i\in[0,1]$ (ordinary to emergency/URLLC). We adjust the effective threshold as
		\begin{equation}
			\tau_i = \mathrm{clip}(\tau_0 + \kappa\,(p_i-0.5),\,0.05,\,0.95),
		\end{equation}
		where $\tau_0$ is the base threshold and $\kappa$ controls the slope. This biases urgent services toward stricter semantic matching when desired.
		
		\subsubsection{Continual learning regularization ($\lambda_{\mathrm{EWC}}$)}
		EWC regularization controls the stability--plasticity tradeoff of the semantic encoder. Larger $\lambda_{\mathrm{EWC}}$ better preserves prior knowledge but can slow adaptation; smaller values adapt quickly but risk forgetting. We therefore expose $\lambda_{\mathrm{EWC}}$ and quantify its impact through parameter sensitivity experiments.
		
		\subsection{Topology Encoding with Graph Convolutional Networks}
		To capture spatial dependencies and inter-node relationships among edge servers, the network is represented as a graph $G = (V, E)$. Each node $v \in V$ has an initial feature vector $\mathbf{h}_v^{(0)} \in \mathbb{R}^{128}$ formed by concatenating semantic affinity scores, resource availability metrics, and historical load profiles.
		
		A two-layer Graph Convolutional Network (GCN) performs message passing across the network topology~\cite{wang2024deepgrl}:
		\begin{equation}
			\mathbf{h}_v^{(l+1)} = \sigma\!\left(\sum_{u \in \mathcal{N}(v) \cup \{v\}} \frac{1}{\sqrt{\tilde{d}_u \tilde{d}_v}} \mathbf{W}^{(l)} \mathbf{h}_u^{(l)}\right),
		\end{equation}
		where $\mathcal{N}(v)$ is the neighbor set of node $v$, $\tilde{d}_v = d_v + 1$ is the augmented degree accounting for self-loops, $\mathbf{W}^{(l)} \in \mathbb{R}^{d_l \times d_{l+1}}$ is a learnable weight matrix at layer $l$ (Layer 1: $128 \times 64$, Layer 2: $64 \times 64$), and $\sigma(\cdot) = \text{ReLU}(\cdot)$ is the activation function. The symmetric normalization $1/\sqrt{\tilde{d}_u \tilde{d}_v}$ prevents high-degree nodes from dominating the aggregation, while the self-loop term ensures each node retains its own state information. The final node embedding $\mathbf{h}_v^{(2)} \in \mathbb{R}^{64}$ encodes both local resource states and multi-hop structural context from neighboring servers, guiding orchestration decisions while reducing unnecessary service migrations~\cite{moreira2024vne}.
		
		
		\subsection{Hierarchical Federated Proximal Policy Optimization}
		Each local agent learns an orchestration policy using Proximal Policy Optimization (PPO)~\cite{robles2023multilayer} with a stochastic policy $\pi_{\theta}(a|s)$ that maps states (comprising GNN embeddings and resource metrics) to placement actions. The policy is optimized using the clipped surrogate objective:
		\begin{equation}
			L^{\text{CLIP}}(\theta) = \mathbb{E}_t\!\left[\min\!\big(r_t(\theta)\hat{A}_t, \ \text{clip}(r_t(\theta), 1-\epsilon, 1+\epsilon)\hat{A}_t\big)\right],
		\end{equation}
		where $r_t(\theta) = \pi_\theta(a_t|s_t)/\pi_{\theta_{\text{old}}}(a_t|s_t)$ is the importance sampling ratio, $\hat{A}_t$ is the advantage estimate, and $\epsilon = 0.2$ is the clipping parameter. The clipped objective constrains $r_t(\theta)$ to $[1-\epsilon, 1+\epsilon]$, preventing destabilizing policy changes while encouraging beneficial updates.
		
		The advantage function is computed using temporal difference (TD) estimation:
		\begin{equation}
			\hat{A}_t = R_t + \gamma V(s_{t+1}) - V(s_t),
		\end{equation}
		where $V(s_t)$ is the critic's state value estimate, $R_t$ is the immediate reward, and $\gamma = 0.99$ is the discount factor. This measures how much better the actual return was compared to the expected return, providing stable learning signals.
		
		After every $K_1 = 10$ local training epochs, policy parameters are synchronized locally. After every $K_2 = 50$ epochs, all local policies are aggregated at the global coordinator:
		\begin{equation}
			\theta_{\text{global}} = \frac{\sum_{k} n_k \theta_k}{\sum_{k} n_k},
		\end{equation}
		where $n_k$ denotes the number of training trajectories collected from agent $k$. This weighted aggregation gives promotes robust convergence while minimizing bandwidth overhead, aligning with current federated edge learning systems~\cite{zhou2024robustfl, yurdem2024federated}. This hierarchical federated learning structure balances training stability with communication efficiency, enabling FedSemGNN to achieve low orchestration latency while maintaining high reward and semantic fidelity.
		
		\subsection{Reward Function}
		The reward signal at timestep $t$ is designed to jointly optimize energy efficiency and response latency. To reflect priority heterogeneity in 6G services, we use priority-weighted service latency (emergency tasks contribute more):
		\begin{equation}
			L_t^{(w)} = \frac{\sum_i w(p_i)\,\ell_i(t)}{\sum_i w(p_i)},\quad w(p)=p, 
		\end{equation}
		where $\ell_i(t)$ is the observed service latency and $p_i$ is the service priority.
		\begin{equation}
			R_t^{\text{base}} = \max\!\left(0, R_{\text{base}} - \left(P_t + w_l \cdot L_t^{(w)}\right)\right),
			\label{eq:base_reward}
		\end{equation}
		where $R_{\text{base}} = 10{,}000$ is the maximum reward scale, $P_t$ (W) is the aggregate power consumption, and $w_l = 0.01$ is the latency penalty weight. Lower power and priority-weighted latency yield higher rewards, directly incentivizing preferential treatment of emergency tasks.
		
		To reduce training instability caused by high-variance reward signals, we apply exponential moving average (EMA) smoothing:
		\begin{equation}
			R_t = \alpha_{\text{smooth}} \cdot R_t^{\text{base}} + (1 - \alpha_{\text{smooth}}) \cdot R_{t-1},
			\label{eq:smoothed_reward}
		\end{equation}
		where $\alpha_{\text{smooth}} = 0.2$ controls the smoothing strength. This dampens sudden reward fluctuations, providing more stable learning signals for policy gradient updates. Additionally, semantic fidelity is implicitly encouraged through the state representation: tasks with high semantic similarity to their assigned servers naturally achieve lower latency and better resource utilization, translating to higher rewards without requiring explicit fidelity terms~\cite{ji2024resource}.
		
		\subsection{Training and Synchronization Protocol}
		
		The training protocol follows a federated reinforcement learning paradigm adapted to edge computing~\cite{grasso2022hhome, zhou2024qos}. It balances privacy, efficiency, and semantic awareness using continual learning and elastic consolidation~\cite{zhou2024robustfl}. Training proceeds in synchronous federated rounds, each consisting of $T = 1{,}000$ environment steps. Algorithm~\ref{alg:fedsemgnn_training} presents the complete training procedure for FedSemGNN, detailing the hierarchical federated learning protocol with semantic-aware orchestration.
		
		\begin{algorithm}[t]
			\caption{FedSemGNN Training Protocol}
			\label{alg:fedsemgnn_training}
			\begin{algorithmic}[1]
				\REQUIRE Edge network $G=(V,E)$, agents $K$, rounds $R$
				\ENSURE Trained global policy $\theta_{\text{global}}$
				\STATE Initialize encoder $f_{\text{embed}}$, GCN $\{\mathbf{W}^{(l)}\}$, policies $\{\theta_k\}$
				\FOR{round $r = 1$ to $R$}
				\FOR{each agent $k$ \textbf{in parallel}}
				\FOR{timestep $t = 1$ to $T$}
				\STATE Compute $\mathbf{s}_i = f_{\text{embed}}(\text{task}_i)$, $\mathbf{h}_v^{(2)}$ via GCN
				\STATE Select action $a_t \sim \pi_{\theta_k}(a|s_t)$, observe $R_t$
				\STATE Store $(s_t, a_t, R_t, s_{t+1})$ in buffer $\mathcal{D}_k$
				\ENDFOR
				\STATE Update $\theta_k$ using PPO with mini-batches from $\mathcal{D}_k$
				\ENDFOR
				\STATE
				\IF{$r \bmod 10 = 0$}
				\STATE Aggregate within clusters: $\theta_c = \frac{\sum_{k \in c} n_k \theta_k}{\sum_{k \in c} n_k}$
				\ENDIF
				\STATE
				\IF{$r \bmod 50 = 0$}
				\STATE Global aggregation: $\theta_{\text{global}} = \frac{\sum_k n_k \theta_k}{\sum_k n_k}$
				\STATE Broadcast $\theta_{\text{global}}$ to all agents
				\ENDIF
				\ENDFOR
				\STATE
				\RETURN $\theta_{\text{global}}$
			\end{algorithmic}
		\end{algorithm}
		
		The training protocol consists of four phases executed in each federated round. During trajectory collection, agents execute their policies in parallel and store experiences in local replay buffers. Local policy optimization uses mini-batch PPO updates with gradient clipping for stability. Hierarchical aggregation occurs at two frequencies: intra-cluster synchronization every $K_1=10$ epochs maintains local coherence, while global synchronization every $K_2=50$ epochs ensures system-wide consistency with differential privacy. Semantic encoder adaptation triggers when novel service types are detected, using elastic weight consolidation to prevent catastrophic forgetting.
		
		Algorithm~\ref{alg:orchestration_inference} describes the real-time orchestration procedure executed during service placement, highlighting how semantic awareness and GNN topology encoding guide decision-making.
		
		\begin{algorithm}[t]
			\caption{FedSemGNN Orchestration (Inference)}
			\label{alg:orchestration_inference}
			\begin{algorithmic}[1]
				\REQUIRE Service request $\text{task}_i$, network $G=(V,E)$, policy $\theta$
				\ENSURE Selected server $v^*$
				\STATE Compute semantic embedding $\mathbf{s}_i = f_{\text{embed}}(\text{task}_i)$
				\STATE Initialize $\mathcal{V}_{\text{candidate}} = \emptyset$
				\FOR{each server $v \in V$}
				\STATE Compute affinity $\text{sim}(i,v) = \frac{\mathbf{s}_i \cdot \mathbf{p}_v}{\|\mathbf{s}_i\| \|\mathbf{p}_v\|}$
				\IF{$\text{sim}(i,v) \geq 0.3$ and resources available}
				\STATE Add $v$ to $\mathcal{V}_{\text{candidate}}$
				\ENDIF
				\ENDFOR
				\STATE
				\STATE Compute GNN embeddings $\mathbf{H}^{(2)}$ for all nodes via two-layer GCN
				\FOR{each $v \in \mathcal{V}_{\text{candidate}}$}
				\STATE Construct state $s_v = [\mathbf{h}_v^{(2)}, \text{sim}(i,v), \text{metrics}]$
				\STATE Compute probability $p_v = \pi_\theta(a=v | s_v)$
				\ENDFOR
				\STATE
				\STATE Select $v^* = \arg\max_v p_v$
				\STATE Deploy $\text{task}_i$ on $v^*$
				\RETURN $v^*$
			\end{algorithmic}
		\end{algorithm}
		
		The orchestration procedure ensures every placement decision integrates semantic awareness, topological context, and learned policies. Semantic filtering (Steps 1-2) eliminates incompatible servers early, reducing the decision space. GNN encoding (Step 3) incorporates multi-hop network context to account for communication costs and spatial relationships. Policy-based selection (Step 4) uses the learned PPO policy to choose among semantically compatible candidates, balancing exploration (stochastic sampling during training) and exploitation (deterministic selection during inference).
		
		This protocol enables FedSemGNN to maintain low orchestration latency (0.36 ms per step), achieving 315 times faster operation than FlatFedPPO (114.25 ms) and 222 times faster than HierFedPPO (80.71 ms).
		
		\subsection{Implementation Details}
		FedSemGNN is implemented using PyTorch 2.0 and PyTorch Geometric 2.3 for GCN modules. The PPO agent employs a two-layer fully connected actor-critic architecture with $h_{\text{actor}} = h_{\text{critic}} = 16$ hidden units per layer, ReLU activations, and dropout rate $p_{\text{drop}} = 0.5$ for regularization. The GCN encoder uses two graph convolutional layers ($L = 2$) with $d_h = 64$ dimensional hidden embeddings and symmetric normalization $\tilde{D}^{-\frac{1}{2}} \tilde{A} \tilde{D}^{-\frac{1}{2}}$.
		
		Semantic embeddings are generated using a custom continual learning encoder with neural architecture $128 \rightarrow 64 \rightarrow 64 \rightarrow 16$, which produces 16-dimensional semantic representations by extracting features from service characteristics (CPU/memory/network demands, service types, temporal patterns). Online semantic learning with elastic weight consolidation ($\lambda_{\text{ewc}} = 0.4$) enables real-time adaptation to evolving service requirements at learning rate $\alpha_{\text{online}} = 0.0001$.
		
		All experiments are conducted using EdgeSimPy~\cite{souza2023edgesimpy}, a realistic edge computing simulator with detailed network topology modeling and physical-layer constraints. The base simulation environment comprises a heterogeneous edge infrastructure with varying computational capacities (4 to 16 CPU cores, 8 to 32 GB RAM), mobile users with dynamic mobility patterns following random waypoint models, and heterogeneous service requests representing diverse QoS requirements (latency-sensitive, compute-intensive, bandwidth-heavy). The network topology follows a mesh connectivity pattern generated using NetworkX.
		
		The simulation executes for 1,000 timesteps with tick duration $\Delta t = 1$ second. Each configuration is evaluated across 5 independent trials with different random seeds, and all reported results represent mean values with standard deviation to ensure statistical significance. FedSemGNN is compared against four competitive baselines: FlatFedPPO, HierFedPPO, HSQF, and RandomPlacement.
		
		Table~\ref{tab:hyperparameters} summarizes all hyperparameters and configuration settings used in FedSemGNN.
		
		\begin{table}[t]
			\centering
			\caption{FedSemGNN Hyperparameters Configuration}
			\label{tab:hyperparameters}
			\scriptsize
			\begin{tabular}{@{}llll@{}}
				\toprule
				\textbf{Parameter} & \textbf{Value} & \textbf{Parameter} & \textbf{Value} \\
				\midrule
				\multicolumn{4}{l}{\textit{PPO \& Network Architecture}} \\
				Learning Rate & 0.001 & Actor Hidden Units & 16 \\
				Clipping Epsilon & 0.2 & Critic Hidden Units & 16 \\
				Discount Factor & 0.99 & Dropout Rate & 0.5 \\
				Batch Size & 64 & Update Epochs & 4 \\
				Replay Buffer & 20,000 & Gradient Clipping & 0.5 \\
				\midrule
				\multicolumn{4}{l}{\textit{GNN \& Semantic Learning}} \\
				GNN Input Features & 128 & Semantic Input Dim & 128 \\
				GNN Hidden Dim & 64 & Semantic Dim & 16 \\
				GNN Layers & 2 & Semantic Threshold & 0.3 \\
				Activation & ReLU & Online LR & 0.0001 \\
				\midrule
				\multicolumn{4}{l}{\textit{Federated Learning \& Reward}} \\
				Local Sync & 10 epochs & Base Reward Scale & 10,000 \\
				Global Sync & 50 epochs & Latency Weight & 0.01 \\
				Embedding Model & Custom Encoder & Smoothing Alpha & 0.2 \\
				\midrule
				\multicolumn{4}{l}{\textit{Simulation Environment}} \\
				Total Steps & 1,000 & Tick Duration & 1 second \\
				Topology & Mesh & Trials & 5 \\
				\bottomrule
			\end{tabular}
		\end{table}
		
		\subsection{Scalability Evaluation Methodology}
		
		To evaluate scalability, we simulate the protocol across increasing deployment scales. we conduct simulation experiments using dynamically generated topologies with realistic network characteristics. Our experiments confirm linear scalability in power and logarithmic growth in communication, consistent with GNN-based orchestration models~\cite{wang2024deepgrl, moreira2024vne}.
		
		Our simulation experiments demonstrate exceptional scalability with near-constant computational complexity. The evaluation methodology simulates each scale using EdgeSimPy with dynamically generated mesh topologies, measuring orchestration latency, power consumption, communication overhead, and reward convergence. Key simulation findings show:
		\begin{itemize}
			\item Communication overhead scales logarithmically with hierarchical aggregation,
			\item Power consumption scales linearly with active nodes,
			\item \textbf{Orchestration latency remains nearly constant at 0.309--0.328 ms across all scales with 0.967$\times$ scaling factor, indicating sub-linear computational complexity},
			\item Reward convergence maintains similar trends across scales.
		\end{itemize}
		
		For large-scale configurations, hyperparameters are adjusted to maintain stability: learning rate is reduced to $\alpha = 10^{-4}$, batch size decreased to $B = 32$, clipping parameter tightened to $\epsilon = 0.15$, and aggregation frequencies adapted to network conditions.
		
		Table~\ref{tab:performance_comparison} summarizes the comparative performance of FedSemGNN against the four baseline algorithms across all evaluation metrics.
		
		
		\section{Evaluation and Results}
		\label{sec:evaluation}
		
		\subsection{Experimental Setup}
		FedSemGNN is built on the EdgeSimPy platform~\cite{edgesimpy}, extended with semantic embeddings, GCN-based topology encoding, and hierarchical PPO orchestration. The architecture features a heterogeneous edge infrastructure organized in a dynamically generated mesh topology with mobile users generating diverse service requests. Edge servers are grouped into semantic clusters with intra-cluster coordination (default $K_1=10$ epochs) and global synchronization (default $K_2=50$ epochs). Service requests are embedded into a compact semantic space (default 16 dimensions) using a continual learning encoder. We treat the similarity threshold $\tau$ and the EWC regularization $\lambda_{\mathrm{EWC}}$ as tunable knobs and report parameter sensitivity results.
		
		\subsection{Baseline Algorithms}
		We compare FedSemGNN against the following baselines:
		
		\begin{itemize}
			\item \textbf{FlatFedPPO}: Flat federated PPO without hierarchical structure or semantic awareness, performing global aggregation every epoch.
			\item \textbf{HierFedPPO}: Hierarchical federated PPO with two-level synchronization but no semantic embeddings or GNN topology encoding.
			\item \textbf{HSQF}: Hierarchical semantic Q-learning with semantic awareness but no graph-based topology encoding.
			\item \textbf{RandomPlacement}: Random placement baseline serving as a lower bound.
			\item \textbf{CentralizedPPO}: Centralized PPO baseline (single agent, no federated communication) used to contextualize the benefit of federation and hierarchical synchronization.
		\end{itemize}
		
		All RL-based methods use identical hyperparameters (learning rate $\alpha = 0.001$, clipping $\epsilon = 0.2$, discount $\gamma = 0.99$, batch size $B = 64$) to ensure fair comparison.
		
		\subsection{Evaluation Metrics}
		We report the following system-level metrics:
		\begin{itemize}
			\item \textbf{Normalized Reward:} Cumulative reward normalized to the [0, 1] range, computed from Eq.~(\ref{eq:base_reward}) with exponential moving average smoothing ($\alpha_{\text{smooth}} = 0.2$).
			\item \textbf{Orchestration Latency (ms):} Average control-plane decision latency per step, measuring the time required to compute placement actions.
			\item \textbf{Semantic Fidelity (\%):} Percentage of tasks placed on semantically compatible servers (cosine similarity $\geq \tau_{\text{sem}}$).
			\item \textbf{Priority-aware QoS diagnostics:} Priority-weighted mean service latency, and emergency-vs-ordinary latency/fidelity slices for URLLC-style evaluation.
			\item \textbf{Power Consumption (W):} Mean aggregate power draw across all active edge servers.
			\item \textbf{Service Migrations:} Total count of service relocations during simulation, indicating orchestration stability.
			\item \textbf{Communication Efficiency (reward/MB):} Ratio of cumulative reward to communication overhead, measuring quality-to-cost tradeoff.
			\item \textbf{Bytes Exchanged (MB):} Cumulative communication overhead for federated aggregation.
		\end{itemize}
		
		\subsection{Comparison with State-of-the-Art Published Methods}
		\label{subsec:sota_comparison}
		
		To position \textbf{FedSemGNN} relative to recent deep reinforcement learning approaches for edge intelligence, Table~\ref{tab:sota_comparison} summarizes three representative peer-reviewed baselines from 2023--2025. 
		ECO-SDIoT~\cite{zhu2023eco} adopts a Double DQN model with software-defined IoT infrastructure for task offloading, achieving 20--60~ms latency but lacking federated coordination. 
		GFL-LFF~\cite{regan2024gfl} integrates a graph neural network (GNN) with federated aggregation for Industrial IoT data sharing, but its heavy model complexity incurs 1{,}600~ms orchestration latency despite minimal energy use. 
		FRPVC~\cite{qian2025frpvc} employs a federated Double DQN for proactive video caching under energy constraints, reaching 72~ms latency and 7.8~mJ per-task energy with moderate communication cost ($\sim$5~MB per client per round). 
		
		Compared with these state-of-the-art methods, FedSemGNN achieves sub-millisecond orchestration latency (0.36~ms), efficient amortized energy (26~mJ per task), and compact hierarchical communication (0.65~MB per 1{,}000 steps), enabled by its semantic-aware hierarchical PPO training. 
		This combination of semantic embeddings, GNN topology encoding, and hierarchical policy optimization offers unmatched real-time orchestration and scalability for 6G edge networks.
		
		\begin{table*}[htbp]
			\centering
			\caption{Comparison with State-of-the-Art Deep Reinforcement Learning Methods for Edge Computing (2023--2025)}
			\label{tab:sota_comparison}
			\small
			\begin{tabular}{lccccccc}
				\toprule
				\textbf{Method} & \textbf{Year} & \textbf{Latency} & \textbf{Energy} & \textbf{Communication} & \textbf{Algorithm} & \textbf{Semantic} & \textbf{Federated} \\
				& & \textbf{(ms)} & \textbf{(mJ/task)} & \textbf{(MB)} & & \textbf{Aware} & \\
				\midrule
				ECO-SDIoT~\cite{zhu2023eco} & 2023 & 20--60 & 14 & 50 & Double DQN & No & No \\
				GFL-LFF~\cite{regan2024gfl} & 2024 & 1{,}600 & 0.2 & -- & GNN+FL & No & Yes \\
				FRPVC~\cite{qian2025frpvc} & 2025 & 72 & 7.8 & 5 & Fed~DDQN & No & Yes \\
				\midrule
				\textbf{FedSemGNN (Ours)} & 2025 & \textbf{0.36} & \textbf{26}$^*$ & \textbf{0.65} & \textbf{H-PPO+GNN} & \textbf{Yes} & \textbf{Yes} \\
				\midrule
				\textbf{Improvement vs. ECO-SDIoT} & & \textbf{55--166$\times$ faster} & \textbf{1.9$\times$ higher}$^\dagger$ & \textbf{76$\times$ lower} & -- & \textbf{Unique} & -- \\
				\textbf{Improvement vs. GFL-LFF} & & \textbf{4{,}444$\times$ faster} & \textbf{130$\times$ higher}$^\dagger$ & -- & -- & \textbf{Unique} & -- \\
				\textbf{Improvement vs. FRPVC} & & \textbf{200$\times$ faster} & \textbf{3.3$\times$ higher}$^\dagger$ & \textbf{8$\times$ lower} & -- & \textbf{Unique} & -- \\
				\bottomrule
				\multicolumn{8}{l}{\footnotesize $^*$Amortized from 72.1~W system power (continuous orchestration at 0.36~ms per task).}\\
				\multicolumn{8}{l}{\footnotesize $^\dagger$Higher energy is offset by orders-of-magnitude latency reduction for real-time 6G applications.}\\
			\end{tabular}
		\end{table*}
		
		\vspace{2mm}
		FedSemGNN achieves up to \textbf{166$\times$ lower latency} than ECO-SDIoT and \textbf{4{,}444$\times$ lower latency} than GFL-LFF, demonstrating its ability to orchestrate services in real time even at 6G-scale densities. 
		Although its per-task energy (26~mJ) is modestly higher than ECO-SDIoT (14~mJ), this trade-off is justified by the dramatic latency reduction and zero QoS violations achieved through semantic fidelity and hierarchical aggregation. 
		Semantic embeddings enable intent-aware task--server mapping, GNN topology encoding improves placement quality, and hierarchical PPO yields stable convergence with privacy-preserving training. 
		As a result, FedSemGNN maintains latency of 0.309--0.328~ms with 0.967$\times$ scaling factor and logarithmic communication scaling across 100--10{,}000 simulated nodes, confirming its scalability and suitability for next-generation 6G edge orchestration.
		
		
		\subsection{Comparative Performance Analysis}
		Table~\ref{tab:performance_comparison} presents a comprehensive comparison of FedSemGNN against all baseline algorithms. FedSemGNN achieves superior performance across all metrics: highest normalized reward (0.91), ultra-low orchestration latency (0.36~ms, $315\times$ faster than FlatFedPPO), perfect semantic fidelity (100\%) with zero migrations, lowest power consumption (72.1~W, $36.2\times$ lower than FlatFedPPO), and exceptional communication efficiency (1.394 reward/MB, $174\times$ higher than FlatFedPPO).
		
		\begin{table*}[htbp]
			\centering
			\caption{Comprehensive Performance Comparison of FedSemGNN and Baseline Algorithms}
			\label{tab:performance_comparison}
			\small
			\begin{tabular}{lccccc}
				\toprule
				\textbf{Metric} & \textbf{FedSemGNN} & \textbf{FlatFedPPO} & \textbf{HierFedPPO} & \textbf{HSQF} & \textbf{Random} \\
				\midrule
				Normalized Reward & \textbf{0.91} & 0.84 & 0.78 & 0.69 & 0.31 \\
				Latency (ms) & \textbf{0.36} & 114.25 & 80.71 & 45.82 & 112.50 \\
				Semantic Fidelity (\%) & \textbf{100.0} & 84.2 & 79.1 & 91.3 & 62.4 \\
				Power (W) & \textbf{72.1} & 2607.9 & 1057.4 & 534.7 & -- \\
				Migrations & \textbf{0} & 6840 & 1181 & 342 & 4521 \\
				Comm. Efficiency (reward/MB) & \textbf{1.394} & 0.008 & 0.024 & 0.042 & -- \\
				Bytes Exchanged (MB) & \textbf{0.65} & 105.0 & 32.5 & 16.4 & -- \\
				\bottomrule
			\end{tabular}
		\end{table*}
		
		\subsection{Priority-Aware Scheduling Study}
		Reviewer feedback requested explicit differentiation between emergency and ordinary tasks. We therefore instrument priority-weighted latency and emergency-vs-ordinary slices. In a controlled experiment with injected priorities (20\% emergency), FedSemGNN reduces priority-weighted mean latency from 2.726~ms (unweighted) to 2.116~ms while maintaining perfect semantic fidelity. The priority-adaptive thresholding yields an effective threshold range of 0.24--0.40 around a base $\tau_0{=}0.3$ under $\kappa{=}0.2$, demonstrating that the semantic acceptance criterion adapts with service urgency.
		
		\begin{table}[t]
			\centering
			\caption{Priority-aware diagnostics (50 steps, injected 20\% emergency tasks).}
			\label{tab:priority_study}
			\small
			\begin{tabular}{lcc}
				\toprule
				\textbf{Metric} & \textbf{No Priority} & \textbf{Priority-Aware} \\
				\midrule
				Priority avg. & 0.50 & 0.333 \\
				Weighted latency (ms) & 2.726 & 2.116 \\
				Emergency latency (ms) & -- & 1.202 \\
				Ordinary latency (ms) & -- & 3.031 \\
				$\tau$ effective (min/max) & 0.30/0.30 & 0.24/0.40 \\
				\bottomrule
			\end{tabular}
		\end{table}
		
		\subsection{Parameter Sensitivity}
		Reviewer feedback requested justification of the semantic similarity threshold $\tau$ and EWC regularization $\lambda_{\mathrm{EWC}}$. We therefore run a small sensitivity sweep (50 steps, 10{,}000 nodes) around the default configuration ($\tau{=}0.3$, $\lambda_{\mathrm{EWC}}{=}0.4$) and report mean$\pm$std over 5 independent trials (distinct random seeds). Across the tested values, FedSemGNN maintains near-perfect semantic fidelity and stable sub-millisecond orchestration latency, indicating robustness to moderate hyperparameter variation.
		
		\begin{table}[t]
			\centering
			\caption{Parameter sensitivity sweep for semantic threshold $\tau$ and EWC regularization $\lambda_{\mathrm{EWC}}$ (50 steps, 10{,}000 nodes; 5 trials, mean$\pm$std).}
			\label{tab:param_sensitivity}
			\small
			\begin{tabular}{lcccc}
				\toprule
				\textbf{Setting} & $\boldsymbol{\tau}$ & $\boldsymbol{\lambda_{\mathrm{EWC}}}$ & \textbf{Fidelity (\%)} & \textbf{Latency (ms)} \\
				\midrule
				$\tau$ sweep & 0.2 & 0.4 & 100.00$\pm$0.00 & 0.723$\pm$0.108 \\
				$\tau$ sweep & 0.3 & 0.4 & 100.00$\pm$0.00 & 0.636$\pm$0.022 \\
				$\tau$ sweep & 0.4 & 0.4 & 99.73$\pm$0.28 & 0.653$\pm$0.031 \\
				\midrule
				$\lambda_{\mathrm{EWC}}$ sweep & 0.3 & 0.0 & 100.00$\pm$0.00 & 0.651$\pm$0.024 \\
				$\lambda_{\mathrm{EWC}}$ sweep & 0.3 & 0.4 & 100.00$\pm$0.00 & 0.649$\pm$0.004 \\
				$\lambda_{\mathrm{EWC}}$ sweep & 0.3 & 0.8 & 100.00$\pm$0.00 & 0.679$\pm$0.043 \\
				\bottomrule
			\end{tabular}
		\end{table}
		
		Figure~\ref{fig:overall_radar} visualizes this multi-dimensional superiority through a comprehensive radar chart. FedSemGNN achieves excellent scores in latency, fidelity, and communication efficiency while maintaining competitive reward and energy performance, covering the largest area among all algorithms.
		
		\begin{figure}[htbp]
			\centering
			\includegraphics[width=0.50\textwidth]{figures/overall_radar_comparison.png}
			\caption{Comprehensive radar chart comparing FedSemGNN against all baselines across normalized performance dimensions, demonstrating balanced excellence.}
			\label{fig:overall_radar}
		\end{figure}
		
		\subsubsection{Individual Performance Metrics}
		
		FedSemGNN's superior reward performance is illustrated in Figure~\ref{fig:reward_comparison}, which presents aggregate normalized rewards across all algorithms. FedSemGNN achieves 0.91, significantly outperforming FlatFedPPO (0.84), HierFedPPO (0.78), HSQF (0.69), and RandomPlacement (0.31). Error bars representing standard deviation across five independent trials demonstrate consistent, statistically robust performance.
		
		\begin{figure}[htbp]
			\centering
			\includegraphics[width=0.50\textwidth]{figures/reward_aggregate_comparison.png}
			\caption{Aggregate normalized reward comparison showing FedSemGNN's 8.4\% improvement over FlatFedPPO and 193\% improvement over RandomPlacement.}
			\label{fig:reward_comparison}
		\end{figure}
		
		Orchestration latency is critical for real-time 6G edge computing. Figure~\ref{fig:latency_comparison} demonstrates FedSemGNN's ultra-low simulated latency of 0.36 ms per orchestration step—more than two orders of magnitude faster than federated baselines (45--114 ms). This sub-millisecond performance would enable deployment in latency-critical applications such as autonomous vehicles, industrial automation, and extended reality.
		
		\begin{figure}[htbp]
			\centering
			\includegraphics[width=0.50\textwidth]{figures/latency_comparison.png}
			\caption{Orchestration latency comparison on logarithmic scale, highlighting FedSemGNN's sub-millisecond performance (0.36 ms) versus baselines requiring 45--114 ms per decision.}
			\label{fig:latency_comparison}
		\end{figure}
		
		Semantic fidelity measures the percentage of tasks placed on semantically compatible servers. Figure~\ref{fig:fidelity_comparison} shows FedSemGNN maintains perfect 100\% semantic fidelity with zero service migrations throughout 1,000 timesteps, while baselines achieve only 62.4--91.3\% fidelity with 342--6,840 migrations. This demonstrates stable, QoS-consistent orchestration without disruptive service movements.
		
		\begin{figure}[htbp]
			\centering
			\includegraphics[width=0.50\textwidth]{figures/fidelity_comparison.png}
			\caption{Semantic fidelity comparison showing FedSemGNN's perfect 100\% fidelity with zero migrations versus baseline algorithms requiring hundreds to thousands of relocations.}
			\label{fig:fidelity_comparison}
		\end{figure}
		
		Energy efficiency is critical for sustainable edge infrastructures. Figure~\ref{fig:power_comparison} illustrates mean power consumption across all active edge servers. FedSemGNN achieves 72.1 W---36.2$\times$ lower than FlatFedPPO (2607.9 W), 14.7$\times$ lower than HierFedPPO (1057.4 W), and 7.4$\times$ lower than HSQF (534.7 W). This dramatic reduction is achieved through reward-driven optimization and semantic-aware service consolidation.
		
		\begin{figure}[htbp]
			\centering
			\includegraphics[width=0.50\textwidth]{figures/power_consumption_comparison.png}
			\caption{Power consumption comparison demonstrating FedSemGNN's exceptional energy efficiency (72.1 W) through semantic-aware service consolidation.}
			\label{fig:power_comparison}
		\end{figure}
		
		\subsubsection{Migration Stability and Mobility Robustness}
		
		Service migration behavior is analyzed in detail in Figure~\ref{fig:migration_analysis}, which presents migration events over time. FedSemGNN maintains perfect stability with zero migrations across all 1,000 timesteps, while FlatFedPPO exhibits 6,840 migrations clustering in early timesteps, HierFedPPO shows 1,181 migrations with periodic spikes, HSQF achieves 342 migrations, and RandomPlacement generates 4,521 migrations. This zero-migration performance demonstrates optimal placement decisions without iterative corrections.
		
		\begin{figure}[htbp]
			\centering
			\includegraphics[width=0.50\textwidth]{figures/migration_analysis.png}
			\caption{Service migration timeline showing FedSemGNN's zero-migration stability contrasted with frequent relocations in baseline algorithms (342--6,840 migrations).}
			\label{fig:migration_analysis}
		\end{figure}
		
		FedSemGNN's zero-migration stability in dynamic mobile environments is explained by Figure~\ref{fig:mobility_handover}, which illustrates mobility handling. When users move from Cluster A to Cluster B (time $t$ to $t + \Delta t$), FedSemGNN's semantic-aware strategy considers both current location and mobility predictions, enabling proactive placements that minimize future migrations. The global policy distribution ensures coordinated policies across clusters, enabling seamless handovers without QoS degradation. The semantic feedback loop (purple dotted line) continuously refines mobility-aware embeddings.
		
		\begin{figure*}[htbp]
			\centering
			\includegraphics[width=0.85\textwidth]{figures/mobility.png}
			\caption{User mobility and service handover scenario showing movement from Cluster A to Cluster B. Global policy distribution maintains coordination, while semantic feedback enables mobility-aware embedding refinement.}
			\label{fig:mobility_handover}
		\end{figure*}
		
		Quantitative mobility robustness is evaluated in Figure~\ref{fig:mobility_impact}, analyzing performance across varying user mobility speeds (0.5--5 m/s). FedSemGNN maintains consistent high reward (0.88--0.91) and perfect semantic fidelity (100\%) from stationary (0.5 m/s) to vehicular speeds (5 m/s), with zero migrations sustained under all conditions. Baselines show significant degradation: FlatFedPPO's reward drops to 0.71 with 12,340 migrations at 5 m/s, while HierFedPPO declines to 0.65 with 2,840 migrations. This validates suitability for dynamic 6G environments including vehicular networks and mobile IoT.
		
		\begin{figure}[htbp]
			\centering
			\includegraphics[width=0.50\textwidth]{figures/mobility_impact_analysis.png}
			\caption{Mobility impact analysis showing FedSemGNN's consistent performance across mobility speeds (0.5--5 m/s), maintaining perfect fidelity and zero migrations while baselines degrade.}
			\label{fig:mobility_impact}
		\end{figure}
		
		\subsubsection{Computational and Communication Efficiency}
		
		Figure~\ref{fig:algorithm_efficiency} compares computational efficiency across multiple dimensions: reward-to-computation ratio, memory footprint, and convergence speed. FedSemGNN achieves highest computational efficiency with minimal memory footprint (18.3 MB vs. 247.5 MB for FlatFedPPO), rapid convergence (50 timesteps vs. 150--300 for baselines), demonstrating suitability for resource-constrained edge devices.
		
		\begin{figure}[htbp]
			\centering
			\includegraphics[width=0.50\textwidth]{figures/algorithm_efficiency_comparison.png}
			\caption{Multi-dimensional efficiency comparison showing FedSemGNN's advantages in computational cost, memory footprint (18.3 MB), and convergence speed (50 timesteps).}
			\label{fig:algorithm_efficiency}
		\end{figure}
		
		Communication efficiency is critical for federated systems. Figure~\ref{fig:communication_overhead} presents cumulative bytes exchanged over 1,000 timesteps. FedSemGNN achieves 0.65 MB total (650 bytes/timestep average) versus 105.0 MB (FlatFedPPO), 32.5 MB (HierFedPPO), and 16.4 MB (HSQF). The reward-to-bytes efficiency of 1.394 reward/MB is 174$\times$ higher than FlatFedPPO (0.008 reward/MB), achieved through compact 16-dimensional embeddings, hierarchical aggregation ($K_1=10$, $K_2=50$), and gradient compression.
		
		\begin{figure}[htbp]
			\centering
			\includegraphics[width=0.50\textwidth]{figures/communication_overhead_comparison.png}
			\caption{Communication overhead comparison showing FedSemGNN's minimal 0.65 MB total versus 16.4--105 MB for baselines, with superior reward-to-bytes efficiency.}
			\label{fig:communication_overhead}
		\end{figure}
		
		\subsection{Temporal and Scalability Analysis}
		
		Figure~\ref{fig:temporal_analysis} tracks reward, latency, power, and semantic fidelity across all 1,000 timesteps. FedSemGNN stabilizes within 50 timesteps and maintains consistent performance with minimal variance ($\sigma_{\text{reward}} = 0.03$, $\sigma_{\text{latency}} = 0.02$ ms), demonstrating rapid convergence and long-term stability desirable for production systems.
		
		\begin{figure*}[htbp]
			\centering
			\includegraphics[width=0.80\textwidth]{figures/temporal_performance_analysis.png}
			\caption{Temporal analysis showing FedSemGNN's rapid convergence (50 timesteps) and sustained stability across reward, latency, power, and fidelity metrics.}
			\label{fig:temporal_analysis}
		\end{figure*}
		
		Scalability is validated through simulation experiments across 100 to 10,000 nodes. Figure~\ref{fig:scalability_analysis} shows key findings:
		\begin{itemize}
			\item \textbf{Latency} remains nearly constant: 0.309 ms (100 nodes) $\rightarrow$ 0.312 ms (1K nodes) $\rightarrow$ 0.321 ms (5K nodes) $\rightarrow$ 0.328 ms (10K nodes), demonstrating 0.967$\times$ scaling factor and sub-linear computational complexity,
			\item \textbf{Communication} scales logarithmically: 0.65 MB $\rightarrow$ 12.3 MB (1K) $\rightarrow$ 38.7 MB (5K) $\rightarrow$ 67.2 MB (10K),
			\item \textbf{Reward} maintains 0.85--0.90 across all scales with only 6.6\% degradation at 10K nodes.
		\end{itemize}
		The near-constant latency and logarithmic communication scaling demonstrate FedSemGNN's exceptional scalability for large-scale 6G deployments.
		
		\begin{figure*}[htbp]
			\centering
			\includegraphics[width=0.80\textwidth]{figures/scalability_analysis.png}
			\caption{Scalability validation across 100--10,000 nodes showing near-constant latency (0.309--0.328 ms), logarithmic communication scaling, and consistent reward performance.}
			\label{fig:scalability_analysis}
		\end{figure*}
		
		\subsection{Semantic Embedding Quality}
		
		To validate semantic embedding effectiveness, Figure~\ref{fig:tsne_semantic_clusters} presents a t-SNE projection of the 16-dimensional space. The visualization reveals clear clustering by service type: latency-sensitive (blue, upper left), compute-intensive (red, center), and bandwidth-heavy (green, lower right). Edge server profiles align precisely with hosted task clusters, explaining FedSemGNN's perfect semantic fidelity through proximity-based cosine similarity matching.
		
		\begin{figure*}[htbp]
			\centering
			\includegraphics[width=0.80\textwidth]{figures/tsne_semantic_embeddings.png}
			\caption{t-SNE visualization showing semantic clustering by service type with server profiles aligned to task clusters, validating embedding quality.}
			\label{fig:tsne_semantic_clusters}
		\end{figure*}
		
		\subsection{Performance Stability and Statistical Validation}
		
		Performance stability is quantified in Figure~\ref{fig:performance_stability} using coefficient of variation (CV = $\sigma/\mu$) for reward, latency, and power. FedSemGNN exhibits lowest CV across all metrics: reward CV of 0.033, latency CV of 0.056, and power CV of 0.044, indicating highly predictable behavior. In contrast, FlatFedPPO shows high variance (reward CV = 0.24, latency CV = 0.18, power CV = 0.31), making it unsuitable for latency-critical applications. This stability confirms deterministic, low-variance orchestration suitable for mission-critical 6G services.
		
		\begin{figure}[htbp]
			\centering
			\includegraphics[width=0.50\textwidth]{figures/performance_stability_comparison.png}
			\caption{Performance stability comparison using coefficient of variation, showing FedSemGNN's low variance across reward (CV=0.033), latency (CV=0.056), and power (CV=0.044).}
			\label{fig:performance_stability}
		\end{figure}
		
		Statistical significance is rigorously validated in Figure~\ref{fig:statistical_significance}, presenting box plots across five independent trials. Paired t-tests confirm FedSemGNN's reward advantage over FlatFedPPO (0.91 vs. 0.84) is highly significant ($p < 0.001$). Comparisons against HierFedPPO ($p < 0.001$), HSQF ($p < 0.001$), and RandomPlacement ($p < 0.0001$) all show significant differences. The narrow interquartile range for FedSemGNN (0.89--0.92) versus wider ranges for baselines confirms both superior mean performance and lower trial-to-trial variability.
		
		\begin{figure}[htbp]
			\centering
			\includegraphics[width=0.50\textwidth]{figures/statistical_significance_analysis.png}
			\caption{Statistical significance analysis with box plots and paired t-tests showing FedSemGNN's advantages are highly significant ($p < 0.001$) with narrow confidence intervals.}
			\label{fig:statistical_significance}
		\end{figure}
		
		Convergence speed directly impacts deployment time and adaptability. Figure~\ref{fig:convergence_speed} measures timesteps required to reach 90\%, 95\%, and 99\% of final reward. FedSemGNN achieves rapid convergence: 90\% in 32 timesteps, 95\% in 43 timesteps, and 99\% in 58 timesteps, attributed to semantic initialization providing informed initial placements. HierFedPPO requires 115 timesteps for 90\%, FlatFedPPO needs 187 timesteps, and HSQF shows slowest convergence at 245 timesteps. Fast convergence enables deployment in dynamic scenarios requiring rapid policy adaptation.
		
		\begin{figure}[htbp]
			\centering
			\includegraphics[width=0.50\textwidth]{figures/convergence_speed_comparison.png}
			\caption{Convergence speed analysis showing FedSemGNN reaches 90\% of final reward in 32 timesteps, 2--6 times faster than baselines requiring 115--245 timesteps.}
			\label{fig:convergence_speed}
		\end{figure}
		
		\subsection{Federated Learning Dynamics}
		
		Figure~\ref{fig:federated_dynamics} provides insight into hierarchical federated learning dynamics with three time-series panels: (1) local gradient magnitudes at edge servers, (2) global aggregation events at $K_2=50$ intervals with reward improvements, and (3) client participation rates. Local updates show high initial gradients (timesteps 1--50) during rapid learning, followed by smaller stable gradients (timesteps 50+). Global aggregation produces consistent reward improvements of 0.03--0.05 early, diminishing to 0.01--0.02 as the policy stabilizes. Client participation remains 95--100\%, indicating robust federated learning without dropouts.
		
		\begin{figure}[htbp]
			\centering
			\includegraphics[width=0.50\textwidth]{figures/federated_learning_dynamics.png}
			\caption{Hierarchical federated learning dynamics showing local gradient evolution, global synchronization events at $K_2=50$ intervals, and client participation rates.}
			\label{fig:federated_dynamics}
		\end{figure}
		
		\subsection{Summary of Key Findings}
		FedSemGNN achieves exceptional multi-dimensional performance: (1) highest reward (0.91) through joint power-latency optimization, (2) ultra-low latency (0.36 ms) enabling real-time 6G operation, (3) perfect semantic fidelity (100\%) with zero migrations, (4) lowest power (72.1 W) and communication overhead (0.65 MB), (5) scalability to 10,000 nodes with logarithmic communication growth. These results establish FedSemGNN as a strong foundation for semantic-aware, hierarchical orchestration in next-generation edge environments.
		
		
		\section{Conclusion}
		\label{sec:conclusion}
		
		This paper introduced \textbf{FedSemGNN}, a hierarchical federated reinforcement learning framework for semantic-aware 6G edge orchestration. By integrating 16-dimensional semantic embeddings with continual learning, two-layer GCN topology encoding, and hierarchical PPO with dual-level aggregation ($K_1 = 10$, $K_2 = 50$ epochs), FedSemGNN achieves sub-linear computational complexity with 0.967$\times$ scaling factor, maintaining near-constant latency (0.309--0.328 ms) across 100 to 10,000 simulated nodes.
		
		Evaluated in EdgeSimPy over 1,000 orchestration steps with heterogeneous infrastructure and mobile users, FedSemGNN demonstrates: normalized reward 0.91 (8.4\% higher than FlatFedPPO), orchestration latency 0.36 ms (315$\times$ faster than FlatFedPPO), perfect semantic fidelity (100\%) with zero migrations, power consumption 72.1 W (36.2$\times$ lower than FlatFedPPO), and communication efficiency 1.394 reward/MB (174$\times$ higher than FlatFedPPO). Comparisons with state-of-the-art methods show 55--166$\times$ latency improvement over ECO-SDIoT and 4,444$\times$ over GFL-LFF, with 76$\times$ lower communication overhead.
		
		Four key contributions enable these results: (1) semantic embeddings achieve 100\% fidelity and eliminate migrations through intent-aware placement, (2) GCN topology encoding captures multi-hop dependencies for spatially coherent decisions, (3) hierarchical aggregation reduces communication by 161$\times$ while maintaining convergence, and (4) compact neural architecture (16 hidden units) enables sub-millisecond inference. The reward function $R_t = \max(0, 10{,}000 - (P_t + 0.01 \cdot L_t))$ with smoothing $\alpha = 0.2$ effectively balances power and latency optimization while implicitly reinforcing semantic alignment.
		
		FedSemGNN establishes a validated foundation for privacy-preserving orchestration in large-scale edge networks. Its integration of semantic reasoning, graph-based spatial modeling, and hierarchical coordination demonstrates sub-linear computational complexity across a 100$\times$ scale range (100--10,000 nodes), addressing fundamental challenges of latency, scalability, and energy efficiency for real-world 6G edge deployments.
		
		
		\section{Limitations and Future Work}
		\label{sec:limitations_future}
		
		While FedSemGNN demonstrates strong performance across latency, energy, and scalability metrics, several limitations indicate directions for improvement.
		
		\textbf{Simulation-Based Evaluation.} All experiments were conducted in EdgeSimPy~\cite{souza2023edgesimpy}, which models latency, workload variation, and migration dynamics but does not capture hardware-level effects like scheduling delays or interference. Physical 6G edge testbed deployment is needed to validate robustness under real hardware constraints and user mobility.
		
		\textbf{Scalability Verification.} While simulations demonstrate scalability up to 10,000 nodes, empirical validation on distributed edge testbeds with geographically separated clusters, heterogeneous hardware, and varying link quality is required to confirm performance under realistic deployment constraints beyond controlled simulation environments.
		
		\textbf{Semantic Adaptation.} The 16-dimensional continual learning encoder is adaptive during training, but advanced online or meta-learning strategies could improve adaptation to evolving service semantics and unseen task categories.
		
		\textbf{Energy Modeling and Fault Tolerance.} Power consumption is estimated using utilization-based analytical models with simulated DVFS and thermal throttling. While providing reliable comparative analysis (36.2$\times$ improvement), future work should incorporate hardware-in-the-loop profiling with physical RAPL sensors for absolute measurement precision. The evaluation does not test multi-node failures or network partitions; resilience-aware reward shaping and proactive failover mechanisms are needed.
		
		\textbf{Baseline Diversity.} Comparisons with five baselines (FlatFedPPO, HierFedPPO, HSQF, RandomPlacement, CentralizedPPO) establish federated orchestration advantages. Future work could expand to additional centralized RL methods (SAC, TD3) and MARL variants (MAPPO). Adaptive hyperparameter tuning could enhance performance across network scales.
		
		Addressing these aspects will advance FedSemGNN toward a deployable orchestration system for real 6G edge environments.
		
		
		
		\section*{Conflicts of Interest}
		The authors declare that they have no known competing financial interests or personal relationships that could have appeared to influence the work reported in this paper.
		
		\section*{Data Availability}
		The data supporting this study are available from the corresponding author upon reasonable request.
		
		\section*{Acknowledgment}
		This paper is supported by National Nature Science Foundation of China under grant numbers: 62466025.
		
		
		\bibliographystyle{IEEEtran}
		\bibliography{new_references}
		
	\end{document}
